\section*{Capítulo \ref{ch:g2:balance-equilibrium} --- Equilibrio: cuando nada se mueve}

\begin{solution}{exo:g2:balance-equilibrium:1}
La Tierra tira de la lámpara hacia abajo; el escritorio la empuja
hacia arriba, exactamente con la misma intensidad. Las dos \omterm{def:g2:pushes-and-pulls:force}{fuerzas}
quedan en tablas --- \omterm{def:g2:balance-equilibrium:equilibrium}{equilibrio} ---, así que la lámpara se queda
donde está.
\end{solution}

\begin{solution}{exo:g2:balance-equilibrium:2}
Cuando los dos equipos tiran con la misma intensidad en sentidos
contrarios: los tirones se anulan. Las tablas se rompen en el momento
en que un equipo tira más fuerte que el otro.
\end{solution}

\begin{solution}{exo:g2:balance-equilibrium:3}
Se queda horizontal --- pesos iguales a distancias iguales. Cuando una
de las gemelas se aleja, baja su lado: la distancia también cuenta.
\end{solution}

\begin{solution}{exo:g2:balance-equilibrium:4}
Cerca del punto de apoyo, con Lina bien alejada en su lado. La regla:
un jinete más pesado cerca del centro puede empatar con un jinete más
ligero bien alejado --- cuentan tanto el peso como la distancia.
\end{solution}

\begin{solution}{exo:g2:balance-equilibrium:5}
El \omterm{def:g2:balance-equilibrium:point}{punto de equilibrio} de la regla está en su centro; el del martillo,
cerca de la cabeza pesada. Comprobación: apoya cada uno sobre un dedo
y busca el sitio en el que se queda horizontal.
\end{solution}

\begin{solution}{exo:g2:balance-equilibrium:6}
Se encuentran cerca del cepillo --- el extremo pesado. El
\omterm{def:g2:balance-equilibrium:point}{punto de equilibrio} se coloca hacia el lado pesado de un objeto.
\end{solution}

\begin{solution}{exo:g2:balance-equilibrium:7}
La pirámide ancha: puede inclinarse muchísimo antes de que su peso
quede fuera de su base ancha. La torre estrecha vuelca a la mínima
inclinación.
\end{solution}

\begin{solution}{exo:g2:balance-equilibrium:8}
La pértiga hace más lento cada bamboleo y le da tiempo a la
funambulista a volver a colocarse sobre la cuerda. En un bordillo tú
abres los brazos por la misma razón: bamboleos más lentos, más fáciles
de atrapar.
\end{solution}

\begin{solution}{exo:g2:balance-equilibrium:9}
Para sostenerte solo sobre el pie derecho, tu cuerpo tiene que
inclinarse hacia la derecha, de modo que tu peso quede encima de ese
pie --- igual que la torre mantiene su peso encima de su base. La
pared está justo donde necesitas inclinarte: te prohíbe el
movimiento, tu peso se queda colgando sobre el hueco vacío del pie
levantado y te vas de lado.
\end{solution}

\begin{solution}{exo:g2:balance-equilibrium:10}
La grúa es un balancín gigante cuyo punto de apoyo es la torre. Jinete
uno: la carga, bien alejada en el largo brazo delantero. Jinete dos:
los bloques de hormigón --- mucho más pesados, así que pueden sentarse
cerca, en el corto brazo trasero. Pesado y cerca empata con ligero y
lejos, justo la regla del balancín; el brazo trasero es corto porque
su jinete es el peso pesado.
\end{solution}
